% !TEX root = ../main.tex
\section{Q10: Individuell prestasjonsmåling og -belønning -- formål og begreper}\label{q:10}

\begin{quote}
\emph{I forbindelse med individuel præstationsmåling og -belønning bedes der redegjort for formålene hermed. I forlængelse heraf redegøres for og diskuteres begreberne: a) ratchet-system / ratchet-effekt, b) gaming, c) absolut / relativ præstationsmåling, d) subjektiv / objektiv præstationsmåling.}
\end{quote}

\subsection*{Svar (kort)}
Individuell prestasjonsmåling tjener fire \emph{simultane} formål -- feedback, koordinasjon, belønning/sanksjon og støtte til den organisatoriske arkitekturen -- og fordi samme måltall presses inn i alle fire roller, oppstår de fire problemene spørsmålet ber om: ratchet-effekten, gaming, valg mellom absolutt og relativ måling, og valg mellom subjektiv og objektiv evaluering. Løsningen i BSZ kap.\ 16 er ikke ``bedre måling'' alene, men strukturelt smartere systemdesign: hybrider, bonusbanker, balansert målekort og Jensens lineære pay-for-performance. [SLIDE]

\subsection*{Formålene med individuell prestasjonsmåling og -belønning}

Se utdypning i \thref{performance-measurement-purposes}{formålene med prestasjonsmåling}. BSZ kap.\ 16 angir fire formål, som ofte drar i ulike retninger og er kilden til de etterfølgende delspørsmålene: [SLIDE]

\begin{enumerate}\setlength{\itemsep}{2pt}
\item \textbf{Feedback (læring).} Den ansatte må vite hvor hun står for å justere atferd. Krever rikt, kontinuerlig signal.
\item \textbf{Koordinasjon.} Felles måltall fungerer som retningsangiver på tvers av enheter -- definerer hva organisasjonen prioriterer.
\item \textbf{Belønning og sanksjon.} Måltallet er input til kompensasjon, forfremmelse og oppsigelse. Krever verifiserbarhet, kontrollerbarhet, informativeness.
\item \textbf{Arkitekturstøtte.} Måltallet binder \SOne\ beslutningsrettigheter og \SThree\ insentiver sammen -- ellers oppstår misalignment og residual loss (jf.\ \thref{agency-costs}{agency costs}).
\end{enumerate}

\textbf{Sentralt poeng:} Et måltall som er optimalt for ett formål er sjelden optimalt for de tre andre. Det er denne flerbruken som genererer ratchet, gaming og hele subjektiv/objektiv-debatten. [SLIDE]

\subsection*{Relevante begreper (kort)}

\paragraph{a) \thref{ratchet-effect}{Ratchet-system og ratchet-effekt}\quad\STwo/\SThree}
Et \emph{ratchet-system} er målsetting som bygger neste periodes mål på inneværende periodes ytelse. \emph{Ratchet-effekten} er agentens rasjonelle motrespons: hold igjen i dag for å unngå strammere mål i morgen -- ekstremvarianten er \emph{sandbagging}. Forutsetter gjentakelse, asymmetrisk informasjon og manglende commitment fra prinsipalen. Historisk dokumentert i sovjetisk planøkonomi (Berliner, Weitzman). Løsning: flerårige LTI, eksterne benchmarks (relativ måling), Jensens lineære pay-for-performance som bryter koblingen mellom budsjett og bonus. [SLIDE/BOK]

\paragraph{b) \thref{gaming-in-performance-measurement}{Gaming og horisontproblem}\quad\STwo}
\emph{Gaming} er manipulasjon av selve målet i stedet for å skape verdien målet representerer; \emph{horisontproblemet} er kortsiktig optimering som ofrer langsiktig verdi. Teoretisk rot: Goodhart's law -- ``når en måling blir et bonusmål, slutter den å være et godt mål.'' Holmström \& Milgrom (1991) viser at sterk insentivvekt på ett målbart mål vrir innsats bort fra ikke-målbare oppgaver (multitasking-problem). Løsninger: balansert målekort, bonusbanker, claw-back, subjektiv overstyring. [SLIDE]

\paragraph{c) \thref{absolute-vs-relative-performance}{Absolutt vs.\ relativ prestasjonsmåling}\quad\STwo}
Absolutt: målt mot fast standard, uavhengig av andre. Relativ: målt mot peer-gruppe (kolleger, eksterne benchmarks). Teorigrunnlag: Holmström (1979) \emph{informativeness principle} -- relativ måling filtrerer ut felles sjokk og gjør målet mer informativt om innsats, slik at insentivkoeffisient \thref{incentive-coefficient}{$\beta$} kan settes høyere uten å overlaste agenten med risiko. Pris: rat race, sabotasje, kunnskapshamstring (Microsoft stack ranking). Hybrid (absolutt budsjett + rTSR-overlegg) er praksis i de fleste børsnoterte LTI-programmer. [SLIDE/BOK]

\paragraph{d) \thref{subjective-vs-objective-performance}{Subjektiv vs.\ objektiv prestasjonsmåling}\quad\STwo}
Objektiv: forhåndsdefinerte, verifiserbare, kvantitative kriterier. Subjektiv: ex post helhetsvurdering av evaluator. Objektive mål er juridisk holdbare og enkle å kommunisere, men dekker bare målbare jobbdimensjoner og er gaming-utsatte. Subjektive evalueringer fanger samarbeid, lederskap og mentoring, men introduserer påvirkningskostnader, partiskhet og supervisor shirking. Nesten alltid hybrid: objektiv formel + subjektiv overstyring, gjerne kalibrert på tvers av evaluatorer. [SLIDE]

\subsection*{Illustrasjon}

\begin{figure}[h]
  \centering
  \includegraphics[width=0.6\linewidth]{BSZ_kap_16/by_slide/slide_016/BSZ_kap_16_slide_016_asset_01_image3.png}
  \caption{Tradisjonelt budget-bonus-system. Den knekkede formen (hurdle ved 80\,\%, cap ved 120\,\%) gir både gaming (rett under hurdle: skyv resultat til neste år; rett over cap: hold igjen) og ratchet-effekt (neste års mål settes fra årets resultat). (BSZ kap.\ 16, slide 16) [SLIDE]}
\end{figure}

\begin{figure}[h]
  \centering
  \includegraphics[width=0.6\linewidth]{BSZ_kap_16/by_slide/slide_019/BSZ_kap_16_slide_019_asset_01_image4.png}
  \caption{Jensens forslag: lineær pay-for-performance der bonus er uavhengig av budsjettet. Budsjettet brukes til planlegging og koordinasjon (formål 2), mens belønning (formål 3) er frikoblet -- da forsvinner motivasjonen for sandbagging og ratchet-tilpasning. (BSZ kap.\ 16, slide 19) [SLIDE]}
\end{figure}

\subsection*{Real-world eksempler}

\textbf{Eks.~1 -- Wells Fargo (US, 2016).}\quad [EKS]
Wells Fargo hadde aggressive cross-sell-kvoter for filialer: bonus knyttet objektivt og absolutt til antall produkter solgt per kunde. Resultatet var 3{,}5 mill.\ falske kundekontoer åpnet uten kundenes viten -- klassisk \emph{gaming} av et énmålssystem. Skandalen utløste \$185 mill.\ bot og styrtende aksjekurs. Etter skandalen omla Wells Fargo til et mer balansert system med subjektiv kvalitetsvurdering, kundetilfredshetsmål og lavere insentivvekt på ren volumvekst. Eksempelet illustrerer hvordan \emph{ren objektiv, absolutt} måling kombinert med høy $\beta$ er en gaming-katastrofe ventende.

\textbf{Eks.~2 -- Microsoft (skift 2013, sammenligning).}\quad [EKS]
Microsoft brukte fra ca.\ 2000 til 2013 ``stack ranking'': en tvungen \emph{relativ} rangering der ledere måtte plassere ansatte i forhåndsdefinerte distribusjoner og bunn-10\,\% ble vurdert ut. Systemet skapte sabotasje (ingen ville dele kunnskap med en konkurrent), kunnskapshamstring og ``hold dårlige folk i teamet for å se bedre ut selv''-strategier. I 2013 avskaffet Steve Ballmer/Satya Nadella stack ranking og innførte ``Connect''-samtaler -- en mer \emph{subjektiv, mindre relativ} evaluering med fokus på vekst og samarbeid. Sammenligningen med Wells Fargo viser at \emph{begge ytterpunktene} kan feile: Wells Fargo feilte med for objektivt og absolutt; Microsoft feilte med for relativt. Begge endte i hybrid.

\subsection*{Kobling til Mærsk Power-to-X}
PtX-prosjektledelse illustrerer flere av problemene konkret: [CASE]
\begin{itemize}\setlength{\itemsep}{2pt}
\item \textbf{Ratchet-fare} hvis ``anlegg-ferdigstilt-per-år'' brukes som mål uten korreksjon for prosjektmodning -- lederen lærer å holde igjen for å unngå strammere mål neste år.
\item \textbf{Gaming/horisontproblem} hvis kortsiktig OPEX brukes som måltall -- vedlikehold på elektrolysører utsettes for å pynte på årets tall.
\item \textbf{Løsning -- hybrid:} Objektive milestones (FID, kontrakter signert, IRR-måltall) for verifiserbarhet og juridisk holdbarhet, kombinert med \emph{subjektiv} ratifisering av kvalitet og partnervalg via \thref{decision-management-control}{Decision Control}-instans i konsernledelsen.
\item \textbf{Relativ benchmarking} mot andre PtX-prosjekter (Ørsted, Air Liquide) filtrerer ut bransje-sjokk (elpris, hydrogen-etterspørsel) og isolerer lederens reelle bidrag (jf.\ \thref{controllability-principle}{controllability}).
\end{itemize}

\subsection*{Fallgruber}
\begin{itemize}\setlength{\itemsep}{2pt}
\item Glemme \emph{formålene} -- spørsmålet ber eksplisitt om dem først. Sensorer venter alle fire (feedback / koordinasjon / belønning / arkitekturstøtte).
\item Behandle de fire begrepene som uavhengige. De er \emph{konsekvenser} av at samme måltall presses i flere formål -- vis sammenhengen.
\item Forveksle ratchet (dynamisk problem på tvers av perioder) med gaming (statisk manipulasjon innen perioden). Sandbagging er en bro: spesifik form for gaming \emph{motivert} av ratchet.
\item Beskrive relativ måling kun positivt -- nevn sabotasje, rat race, peer-pool manipulasjon.
\item Påstå at subjektiv = irrasjonell. Subjektiv evaluering er ofte mer informativ enn dårlige objektive måltall.
\item Mangle hybridløsningen. ``Verken eller, men begge med rett vekting'' er ofte BSZ-svaret.
\item Glemme Jensen-kritikken og lineær pay-for-performance som strukturell løsning på ratchet-problemet.
\item Glemme balansert målekort som svar på multitasking/gaming.
\end{itemize}

\subsection*{Håndskrivbar disposisjon (15-min forberedelse)}
\begin{itemize}\setlength{\itemsep}{2pt}
\item \textbf{Tese:} Måling tjener \emph{fire} formål samtidig -- konfliktene mellom dem genererer alle problemene i a)--d). Løsning: smartere systemdesign, ikke bare ``bedre måling''.
\item \textbf{Fire formål} (\thref{performance-measurement-purposes}{purposes}): (i) Feedback, (ii) Koordinasjon, (iii) Belønning/sanksjon, (iv) Arkitekturstøtte. Trade-offs mellom dem $\Rightarrow$ resten av spørsmålet.
\item \textbf{a) \thref{ratchet-effect}{Ratchet}:} Mål fra historie $\Rightarrow$ sandbagging. Sovjet-eksempel. Løsning: flerårig LTI, ekstern benchmark, lineær pay (Jensen).
\item \textbf{b) \thref{gaming-in-performance-measurement}{Gaming + horisont}:} Goodhart, Holmström \& Milgrom multitasking. Løsning: BSC, bonusbank, claw-back, subj.\ overstyring.
\item \textbf{c) \thref{absolute-vs-relative-performance}{Absolutt vs.\ relativ}:} Holmström '79 informativeness $\Rightarrow$ relativ filtrerer felles støy, lavere risikopremie, men sabotasje. Hybrid (rTSR).
\item \textbf{d) \thref{subjective-vs-objective-performance}{Subj.\ vs.\ obj.}:} Obj.\ verifiserbar men gaming-utsatt; subj.\ fanger multidim.\ men gir påvirkningskost + bias. Hybrid.
\item \textbf{Illustrasjon:} BSZ kap.\ 16 slide 16 (knekket budget-bonus) vs.\ slide 19 (Jensens lineære).
\item \textbf{Eks.\ 1 -- Wells Fargo 2016:} objektiv + absolutt + høy $\beta$ $\Rightarrow$ 3,5 mill.\ falske kontoer.
\item \textbf{Eks.\ 2 -- Microsoft pre/post-2013:} stack ranking (relativ + objektiv) $\Rightarrow$ sabotasje; ``Connect'' (mer subj.\ + absolutt).
\item \textbf{Case:} Mærsk PtX -- milestone-mål (obj.) + subj.\ kvalitetsratifisering + relativ benchmark mot Ørsted/Air Liquide.
\item \textbf{Søyle:} \STwo\ primært; kobler til \SOne\ (controllability) og \SThree\ ($\beta$-design).
\end{itemize}
